\multicolumn{4}{@{}l}{\textbf{chromdata.zarr}}\\*[2pt]
spot coordinates (x, y, z) & \mNative{} native & coords (n\_spots, 3) + layers; unit in uns (ChromData format specification) & dedicated slot: coords (n\_spots, 3) \\
cell › trace › spot hierarchy & \mNative{} native & spots → traces → cells via trace\_id / cell\_id; traces and cells tables (ChromData format specification) & dedicated slot: spots.trace\_id / cell\_id → traces / cells \\
per-spot signals & \mNative{} native & spot\_tracks, row-aligned with spots (ChromData format specification) & dedicated slot: spot\_tracks \\
per-locus (bin) signals & \mNative{} native & bins (locus axis, bin\_id) + bin\_tracks / binm (ChromData format specification) & dedicated slot: bins + bin\_tracks \\
cell table & \mNative{} native & cells table indexed by cell\_id (ChromData format specification) & dedicated slot: cells table \\
cell embeddings (multi-dimensional) & \mNative{} native & cellm: per-cell multi-dimensional arrays (embeddings, UMAP) (ChromData format specification) & dedicated slot: cellm['X\_umap'] (n\_cells, 2) \\
cell spatial position (centroid / outline) & \mNative{} native & the cell spatial position schema of the format specification: cells columns centroid\_x / centroid\_y [/ centroid\_z] described by uns['cell\_spatial'] (unit, frame fov / tissue / global, region, measured / inferred + source, voxel size); outlines as cell\_shapes, GeoParquet WKB polygons (tables/cell\_shapes/<key>.parquet, the SpatialData shapes representation); set\_cell\_positions / cell\_positions / set\_cell\_shapes; loaders fill it (seqFISH+, FOF-CT cell and mapping tables) (ChromData format specification) & dedicated slot: set\_cell\_positions: cells centroid\_x / y / z + uns['cell\_spatial'] record (unit, frame, status); set\_cell\_shapes: outlines as GeoParquet tables/cell\_shapes/cell.parquet (format specification, cell spatial position) \\
contact maps & \mLinked{} linked (external file, queried through the index) & linked, not copied: uns['linked\_scool' / 'linked\_cool'] records path, bin size and cell map; load\_linked\_scool and the group pseudo-bulk query read them (test: a real .scool read back after the round trip). Our own call: the pixels stay in the external .scool / .cool by design. (ChromData format specification) & external file: link\_scool → uns['linked\_scool'] record; load\_linked\_scool(cell='c1') reads the external .scool after the round trip \\
genomic intervals (TADs / loops) & \mNative{} native & intervals: typed IntervalTable (domain / pair / peak / segment); TADs as domain, loops as pair (ChromData format specification) & dedicated slot: intervals: IntervalTable kind='domain' (TADs) and kind='pair' (loops) \\
non-genomic points (RNA spots) & \mNative{} native & points: non-genomic 3-D points in the coords frame (ChromData format specification) & dedicated slot: points['rna'] \\
analysis results with provenance & \mNative{} native & results: ResultsStore of ResultRecord(kind, params, function, version, inputs, created\_utc) (ChromData format specification) & dedicated slot: results: ResultRecord(function, params, ...) \\
partial reads by cell / chromosome & \mNative{} native & read(backed=True): get\_cell / get\_trace / get\_chrom read only the needed row ranges (index/ runs; Fig. 2d) (ChromData format specification) & dedicated slot: read(backed=True).get\_cell / get\_chrom (index/ row ranges) \\
\midrule
\multicolumn{4}{@{}l}{\textbf{FOF-CT (4DN)}}\\*[2pt]
spot coordinates (x, y, z) & \mNative{} native & core table X, Y, Z (required) + \#\#XYZ\_Unit (\url{https://fish-omics-format.readthedocs.io/en/latest/core.html}) & specification only \\
cell › trace › spot hierarchy & \mNative{} native & Spot\_ID → Trace\_ID → Cell\_ID in the core table; Trace Data and Cell Data tables indexed by Trace\_ID / Cell\_ID (\url{https://fish-omics-format.readthedocs.io/en/latest/core.html}, \url{https://fish-omics-format.readthedocs.io/en/latest/trace.html}, \url{https://fish-omics-format.readthedocs.io/en/latest/cell.html}) & specification only \\
per-spot signals & \mNative{} native & Spot Biological Data / Spot Quality tables indexed by Spot\_ID (user-defined columns declared with \#\textasciicircum{} lines) (\url{https://fish-omics-format.readthedocs.io/en/latest/bio.html}, \url{https://fish-omics-format.readthedocs.io/en/latest/quality.html}) & specification only \\
per-locus (bin) signals & \mConv{} by convention & none of the 12 tables is indexed by genomic locus; a per-locus value can only be repeated on every spot as an optional column (\url{https://fish-omics-format.readthedocs.io/en/latest/format_overview.html}) & specification only \\
cell table & \mNative{} native & Cell Data table (4dn\_FOF-CT\_cell), mandatorily indexed by Cell\_ID, with optional user-defined columns (cell type, volume, ...) (\url{https://fish-omics-format.readthedocs.io/en/latest/cell.html}) & specification only \\
cell embeddings (multi-dimensional) & \mConv{} by convention & no embedding slot: the Cell Data table's columns are scalar optional columns, so an embedding is one \#\textasciicircum{} column per dimension (\url{https://fish-omics-format.readthedocs.io/en/latest/cell.html}) & specification only \\
cell spatial position (centroid / outline) & \mNative{} native & Cell/ROI Mapping table (4dn\_FOF-CT\_mapping; conditionally required whenever a Cell Data table is deposited): ROI\_Boundaries of each Cell\_ID in global coordinates, format declared in \#\#ROI\_Boundaries\_Format (OME polygon or OBJ mesh) + \#\#XYZ\_Unit. Outline native; a centroid only as an optional column (\url{https://fish-omics-format.readthedocs.io/en/latest/mapping.html}, \url{https://fish-omics-format.readthedocs.io/en/latest/cell.html}) & specification only \\
contact maps & \mNone{} not supported & no contact-matrix table in the 12 FOF-CT tables (\url{https://fish-omics-format.readthedocs.io/en/latest/format_overview.html}) & specification only \\
genomic intervals (TADs / loops) & \mConv{} by convention & no genomic-interval table (the Sub-Cell ROI table's 'Chromosome\_Domain' is a spatial ROI, not a genomic interval); domain labels only as an optional per-spot column, loops (pairs) not representable (\url{https://fish-omics-format.readthedocs.io/en/latest/format_overview.html}, \url{https://fish-omics-format.readthedocs.io/en/latest/subcell.html}) & specification only \\
non-genomic points (RNA spots) & \mNative{} native & RNA Spot Data table (RNA\_Spot\_ID, X, Y, Z, Gene\_ID, RNA\_name, link ids) (\url{https://fish-omics-format.readthedocs.io/en/latest/rna.html}) & specification only \\
analysis results with provenance & \mConv{} by convention & header \#Software\_* fields describe the tools that produced a table; analysis results only as optional columns (\url{https://fish-omics-format.readthedocs.io/en/latest/format_overview.html}) & specification only \\
partial reads by cell / chromosome & \mConv{} by convention & plain CSV without an index: one cell or chromosome is found only by scanning every row (streamed, e.g. pandas chunksize, so memory stays bounded) (\url{https://fish-omics-format.readthedocs.io/en/latest/format_overview.html}) & specification only \\
\midrule
\multicolumn{4}{@{}l}{\textbf{AnnData}}\\*[2pt]
spot coordinates (x, y, z) & \mConv{} by convention & obsm['spatial'] (n\_obs, 3), a community convention; no trace / unit structure (\url{https://anndata.readthedocs.io/en/stable/generated/anndata.AnnData.html}) & general-purpose slot: spots as obs, obsm['spatial'] (n\_obs, 3) \\
cell › trace › spot hierarchy & \mConv{} by convention & one obs axis: cell\_id / trace\_id only as obs columns; spots and cells cannot both be observations (\url{https://anndata.readthedocs.io/en/stable/generated/anndata.AnnData.html}) & general-purpose slot: trace\_id / cell\_id obs columns \\
per-spot signals & \mNative{} native & spots as obs: X / layers / obs columns. Rated in the most favourable orientation; one AnnData has a single obs axis, so spots, cells and loci cannot all be first-class in the same object (see the hierarchy column). (\url{https://anndata.readthedocs.io/en/stable/generated/anndata.AnnData.html}) & dedicated slot: spots as obs, X \\
per-locus (bin) signals & \mNative{} native & loci as var (cells × loci, as for scATAC peaks): var table. Rated in the most favourable orientation; one AnnData has a single obs axis, so spots, cells and loci cannot all be first-class in the same object (see the hierarchy column). (\url{https://anndata.readthedocs.io/en/stable/generated/anndata.AnnData.html}) & dedicated slot: loci as var, var column \\
cell table & \mNative{} native & cells as obs: obs table. Rated in the most favourable orientation; one AnnData has a single obs axis, so spots, cells and loci cannot all be first-class in the same object (see the hierarchy column). (\url{https://anndata.readthedocs.io/en/stable/generated/anndata.AnnData.html}) & dedicated slot: cells as obs, obs table \\
cell embeddings (multi-dimensional) & \mNative{} native & cells as obs: obsm (X\_umap, X\_pca), the slot designed for per-observation multi-dimensional annotations. Rated in the most favourable orientation; one AnnData has a single obs axis, so spots, cells and loci cannot all be first-class in the same object (see the hierarchy column). (\url{https://anndata.readthedocs.io/en/stable/generated/anndata.AnnData.html}) & dedicated slot: cells as obs, obsm['X\_umap'] \\
cell spatial position (centroid / outline) & \mConv{} by convention & cells as obs: obsm['spatial'] (n\_cells, 2), the squidpy / scanpy convention; not an AnnData schema element (no coordinate system, unit or outline) (\url{https://anndata.readthedocs.io/en/stable/generated/anndata.AnnData.html}, \url{https://squidpy.readthedocs.io/en/stable/notebooks/tutorials/tutorial_visium_hne.html}) & general-purpose slot: cells as obs, obsm['spatial'] (squidpy convention; no coordinate system, no outline) \\
contact maps & \mConv{} by convention & no contact-map type: one dense or sparse locus × locus matrix can be put in varp (loci as var), without pixel/balancing schema, resolutions or per-cell maps. Rated in the most favourable orientation; one AnnData has a single obs axis, so spots, cells and loci cannot all be first-class in the same object (see the hierarchy column). (\url{https://anndata.readthedocs.io/en/stable/generated/anndata.AnnData.html}) & general-purpose slot: one locus × locus matrix in varp \\
genomic intervals (TADs / loops) & \mConv{} by convention & a DataFrame in uns (\url{https://anndata.readthedocs.io/en/stable/generated/anndata.AnnData.html}) & general-purpose slot: DataFrame in uns \\
non-genomic points (RNA spots) & \mConv{} by convention & a DataFrame in uns (or a second AnnData) (\url{https://anndata.readthedocs.io/en/stable/generated/anndata.AnnData.html}) & general-purpose slot: DataFrame in uns \\
analysis results with provenance & \mConv{} by convention & free-form dicts in uns (scanpy writes uns['<tool>']['params']) (\url{https://anndata.readthedocs.io/en/stable/generated/anndata.AnnData.html}) & general-purpose slot: dict in uns \\
partial reads by cell / chromosome & \mNative{} native & Judgement call: anndata.experimental.read\_lazy (0.12, experimental API) opens obs as a lazy xarray Dataset and obsm / X as dask arrays; a cell mask reads the cell\_id column, to\_memory() then reads only the selected chunks. No per-cell index. (\url{https://anndata.readthedocs.io/en/stable/generated/anndata.experimental.read_lazy.html}) & dedicated slot: read\_lazy: obs / obsm / X are lazy (xarray / dask); mask from the cell\_id column, to\_memory() reads only the selected chunks \\
\midrule
\multicolumn{4}{@{}l}{\textbf{SpatialData}}\\*[2pt]
spot coordinates (x, y, z) & \mNative{} native & points element with x, y, z (3-D points are supported) and coordinate-system transformations (\url{https://spatialdata.scverse.org/en/stable/design_doc.html}) & dedicated slot: 3-D points element (x, y, z) \\
cell › trace › spot hierarchy & \mConv{} by convention & cell\_id / trace\_id as columns of the points; a table annotates one element level (region + instance\_key), no nested cell › trace › spot (\url{https://spatialdata.scverse.org/en/stable/design_doc.html}) & general-purpose slot: trace\_id / cell\_id point columns \\
per-spot signals & \mNative{} native & a table annotating the spots points element (region + instance\_key = spot\_id): get\_element\_annotators lists it and get\_values joins it to the points (test); extra columns of the points DataFrame also work (\url{https://spatialdata.scverse.org/en/stable/design_doc.html}) & dedicated slot: a table annotating the spots points (region='spots', instance\_key='spot\_id'); get\_element\_annotators lists it and get\_values joins it to the points \\
per-locus (bin) signals & \mNative{} native & Judgement call: no genomic locus axis among the spatial elements, but a table is an AnnData: loci as var of a cells × loci table annotating the cell shapes, var columns = per-locus signals, as for AnnData (test uses this orientation instead of a dummy table annotating no element). Rated in the most favourable orientation; one AnnData has a single obs axis, so spots, cells and loci cannot all be first-class in the same object (see the hierarchy column). (\url{https://spatialdata.scverse.org/en/stable/design_doc.html}) & dedicated slot: var column of a cells × loci table annotating the cell shapes \\
cell table & \mNative{} native & shapes (or labels) for cells + a table annotating them (\url{https://spatialdata.scverse.org/en/stable/design_doc.html}) & dedicated slot: table annotating the cell shapes \\
cell embeddings (multi-dimensional) & \mNative{} native & obsm of the table annotating the cells (\url{https://spatialdata.scverse.org/en/stable/design_doc.html}) & dedicated slot: obsm['X\_umap'] of the cell table \\
cell spatial position (centroid / outline) & \mNative{} native & shapes element: circles (centroid + radius) or polygons (outline), with coordinate-system transformations (\url{https://spatialdata.scverse.org/en/stable/design_doc.html}) & dedicated slot: shapes elements: circles (centroid + radius) and polygons (outline), each with a coordinate-system transformation \\
contact maps & \mConv{} by convention & no contact-map element: a table is an AnnData, so one locus × locus matrix in varp, as for AnnData (\url{https://spatialdata.scverse.org/en/stable/design_doc.html}) & general-purpose slot: one locus × locus matrix in varp of a table \\
genomic intervals (TADs / loops) & \mConv{} by convention & interval rows as obs of a table annotating no element (\url{https://spatialdata.scverse.org/en/stable/design_doc.html}) & general-purpose slot: interval rows as obs of a table annotating no element \\
non-genomic points (RNA spots) & \mNative{} native & a second points element (feature\_key = gene) (\url{https://spatialdata.scverse.org/en/stable/design_doc.html}) & dedicated slot: points element, feature\_key='gene' \\
analysis results with provenance & \mConv{} by convention & SpatialData.attrs / table uns (free-form) (\url{https://spatialdata.scverse.org/en/stable/design_doc.html}) & general-purpose slot: SpatialData.attrs \\
partial reads by cell / chromosome & \mConv{} by convention & points are lazy dask DataFrames without an index on cell or chromosome: a filter on cell\_id then compute() streams every partition (Fig. 2d, 0.75 s at 1e7 spots). Spatial bounding-box / polygon queries are native but select by space, not by cell or chromosome. (\url{https://spatialdata.scverse.org/en/stable/design_doc.html}) & general-purpose slot: lazy dask points: a filter on cell\_id streams every partition (no index) \\
\midrule
\multicolumn{4}{@{}l}{\textbf{MuData}}\\*[2pt]
spot coordinates (x, y, z) & \mConv{} by convention & obsm['spatial'] of a spots modality (\url{https://mudata.readthedocs.io/stable/io/mudata.html}) & general-purpose slot: obsm['spatial'] of a spots modality \\
cell › trace › spot hierarchy & \mConv{} by convention & spots and cells as separate modalities; the link is an obs column, no parent / child relation (\url{https://mudata.readthedocs.io/stable/io/mudata.html}) & general-purpose slot: trace\_id / cell\_id obs columns of the spots modality \\
per-spot signals & \mNative{} native & a spots modality (spots as obs): X / obs (\url{https://mudata.readthedocs.io/stable/io/mudata.html}) & dedicated slot: spots modality X \\
per-locus (bin) signals & \mNative{} native & var of a cells × loci modality (loci as var). A MuData can hold a spots and a cells modality side by side, but within one modality the orientations exclude each other, as for AnnData (\url{https://mudata.readthedocs.io/stable/io/mudata.html}) & dedicated slot: var of a cells × loci modality \\
cell table & \mNative{} native & global obs or the obs of a cells modality (\url{https://mudata.readthedocs.io/stable/io/mudata.html}) & dedicated slot: obs of a cells modality \\
cell embeddings (multi-dimensional) & \mNative{} native & obsm of a cells modality, or the global obsm (multimodal embeddings) (\url{https://mudata.readthedocs.io/stable/io/mudata.html}) & dedicated slot: obsm['X\_umap'] of a cells modality \\
cell spatial position (centroid / outline) & \mConv{} by convention & obsm['spatial'] of a cells modality (the AnnData / squidpy convention) (\url{https://mudata.readthedocs.io/stable/io/mudata.html}, \url{https://squidpy.readthedocs.io/en/stable/notebooks/tutorials/tutorial_visium_hne.html}) & general-purpose slot: obsm['spatial'] of a cells modality (convention) \\
contact maps & \mConv{} by convention & no contact-map type: varp of a cells × loci modality (one matrix, as for AnnData) (\url{https://mudata.readthedocs.io/stable/io/mudata.html}) & general-purpose slot: one locus × locus matrix in varp \\
genomic intervals (TADs / loops) & \mConv{} by convention & a DataFrame in uns (\url{https://mudata.readthedocs.io/stable/io/mudata.html}) & general-purpose slot: DataFrame in uns \\
non-genomic points (RNA spots) & \mConv{} by convention & an RNA-spot modality with obsm['spatial'] (\url{https://mudata.readthedocs.io/stable/io/mudata.html}) & general-purpose slot: an RNA-spot modality, positions in obsm['spatial'] \\
analysis results with provenance & \mConv{} by convention & free-form uns (\url{https://mudata.readthedocs.io/stable/io/mudata.html}) & general-purpose slot: dict in uns \\
partial reads by cell / chromosome & \mConv{} by convention & read\_h5mu(backed=True) keeps only each modality's X on disk (sliced by row); obs and obsm, i.e. the ids and the coordinates, are loaded whole (test). MuData 0.3 has no lazy reader. (\url{https://mudata.readthedocs.io/stable/io/mudata.html}) & general-purpose slot: read\_h5mu(backed=True): X is sliced from disk by row, but obs and obsm (the coordinates) are loaded whole \\
\midrule
\multicolumn{4}{@{}l}{\textbf{cooler (.cool / .scool)}}\\*[2pt]
spot coordinates (x, y, z) & \mConv{} by convention & extra bin columns x, y, z: one structure per file (or per cell of a .scool, one column set per trace); no spot table (\url{https://cooler.readthedocs.io/en/latest/schema.html}) & general-purpose slot: extra bin columns x, y, z (one structure per file) \\
cell › trace › spot hierarchy & \mNone{} not supported & no trace or spot level: a .scool has per-cell groups, but traces and spots could only be forced in as sets of per-trace bins columns (x\_<trace>, ...; the test shows this writes), which is not a use of the format (rule 6) (\url{https://cooler.readthedocs.io/en/latest/schema.html}) & not counted: cells = .scool cell groups (each with its own bins table); traces = sets of extra bins columns x\_<trace>, y\_<trace>, z\_<trace>; spots = bins rows (writes; not counted, rule 6: not a use of the format) \\
per-spot signals & \mNone{} not supported & no spot level; per-trace intensity columns in each cell's bins table would write (test) but are not a use of the format (rule 6) (\url{https://cooler.readthedocs.io/en/latest/schema.html}) & not counted: one extra bins column per trace (intensity\_<trace>) in each cell's bins table (writes; not counted, rule 6) \\
per-locus (bin) signals & \mNative{} native & bins table, which the schema allows to carry extra columns (e.g. weight, signals) (\url{https://cooler.readthedocs.io/en/latest/schema.html}) & dedicated slot: bins table column \\
cell table & \mConv{} by convention & no cell table, but every cell of a .scool is a cooler with its own metadata attribute: the cell's properties as scalar key-values there (cooler.create\_cooler on /cells/<name>, test); the table is rebuilt by opening every cell group (\url{https://cooler.readthedocs.io/en/latest/schema.html}) & general-purpose slot: each cell's properties as scalar key-values in the metadata attribute of its own cooler (/cells/<name>), written with cooler.create\_cooler(uri, ..., metadata=); a scan of every cell group rebuilds the table \\
cell embeddings (multi-dimensional) & \mNone{} not supported & no per-cell arrays; one scalar metadata key per dimension (umap\_1, umap\_2, ...) writes (test) but is not a use of the format for a 10- to 64-dimensional embedding (rule 6) (\url{https://cooler.readthedocs.io/en/latest/schema.html}) & not counted: one scalar key per embedding dimension (umap\_1, umap\_2) in each cell's metadata (writes; not counted, rule 6) \\
cell spatial position (centroid / outline) & \mConv{} by convention & centroid\_x / centroid\_y (+ a unit key) in each cell's metadata (test); no coordinate frame, no outline (\url{https://cooler.readthedocs.io/en/latest/schema.html}) & general-purpose slot: centroid\_x / centroid\_y (+ a unit key) in each cell's metadata; no frame, no outline \\
contact maps & \mNative{} native & pixels (bulk .cool; per cell .scool; multi-resolution .mcool) (\url{https://cooler.readthedocs.io/en/latest/schema.html}) & dedicated slot: pixels \\
genomic intervals (TADs / loops) & \mConv{} by convention & domain membership as a per-bin label column of the bins table (data, not JSON); loops (pairs) have no place (\url{https://cooler.readthedocs.io/en/latest/schema.html}) & general-purpose slot: per-bin domain label column \\
non-genomic points (RNA spots) & \mNone{} not supported & no place for non-genomic points; JSON text in the metadata attribute works for 2e5 spots (test) but is not counted (\url{https://cooler.readthedocs.io/en/latest/schema.html}) & not counted: 2e5 RNA spots as JSON text in the metadata attribute (stored; not counted, rule 4) \\
analysis results with provenance & \mConv{} by convention & the result as data (a bins column, as for intervals) and its parameters in the free-form JSON metadata attribute next to generated-by / creation-date (provenance metadata counts under rule 4; a result table as JSON would not) (\url{https://cooler.readthedocs.io/en/latest/schema.html}) & general-purpose slot: result as a bins column, its parameters as JSON in the metadata attribute \\
partial reads by cell / chromosome & \mNative{} native & indexed fetch by genomic region (bin1\_offset / chrom\_offset indexes); one cell group of a .scool (\url{https://cooler.readthedocs.io/en/latest/schema.html}) & dedicated slot: indexed fetch('chr2'); one cell group of a .scool \\
(extra test: cell table as JSON) & --- & --- & not counted: a whole cell table as JSON text in the file-level metadata attribute (stored; not counted, rule 4) \\
(extra test: embedding as JSON) & --- & --- & not counted: an embedding matrix as JSON text in the file-level metadata attribute (stored; not counted, rule 4) \\
